\section{Vòng v8: đưa dự báo GNN vào lựa chọn tuyến}\label{sec:m8}
Vòng v8 kiểm tra trực tiếp liệu chất lượng liên kết dự báo có tăng PDR so với
chất lượng hiện tại khi giữ nguyên luật ghim tuyến. Mô hình ở vòng này là
\textbf{Edge-SAGE hồi quy mới}, dự báo tỷ lệ nhận probe trung bình của năm giây
tiếp theo; không phải checkpoint phân loại QoS/survival của Chương~\ref{chap:gnn}.
Mô hình phân loại cũ và bộ hồi quy mới trả lời hai bài toán khác nhau.
GNN ở đây được học có giám sát; lựa chọn tuyến và vị trí theo luật cố định.
V8/v9 không sử dụng học tăng cường hoặc học chính sách hành động từ phần thưởng.
Hình~\ref{fig:v8v9design} tóm tắt yếu tố giữ nguyên và yếu tố thay đổi trong
hai nghiên cứu, giúp tách tác dụng của dự báo đối với mỗi loại hành động.
\Fig{fig_v8_v9_design}{Hai phép so sánh vòng kín: thay thông tin chọn tuyến ở v8, thay trigger chuyển động ở v9. Sơ đồ mô tả thiết kế, không biểu diễn kết quả hiệu quả.}{fig:v8v9design}

\subsection{Dữ liệu phát triển và mô hình}
Chỉ sử dụng trace HOLD đã có của v7: 48 episode test cũ để fit, 20 episode
validation cũ để chọn epoch. Những episode này đã được quan sát ở v7, nên ở
v8 chúng được công bố là \emph{dữ liệu phát triển}. Validation và test v8 dùng
các cặp seed mới, được kiểm toán không trùng các vòng trước. Không chia ngẫu
nhiên các frame của cùng một episode sang cả tập fit và tập chọn epoch.

Tại giây $t$, mỗi nút có tám đặc trưng: tọa độ trừ tâm đội hình (chia 240 m),
vận tốc (chia 10 m/s), bậc hàng xóm và tỷ lệ nhận probe trung bình, chuẩn hóa
theo $N-1$. Mỗi cặp nút có bảy đặc trưng: khoảng cách, RSSI, tỷ lệ nhận hiện
tại, EWMA với hệ số 0,5, chênh tỷ lệ nhận so với giây trước, vận tốc tách xa
và độ lớn vận tốc tương đối. RSSI thiếu được mã hóa thành $-100$ dBm.
Chỉ có cạnh truyền thông điệp khi tỷ lệ nhận hiện tại ít nhất 0,1.
Mọi đầu vào được tính từ quan sát đến giây $t$; không dùng quỹ đạo nhiệm vụ
hoặc gói nhận trong tương lai. Nhãn huấn luyện là
\begin{equation}
 y_{ij,t}=\frac{1}{5}\sum_{h=1}^{5}q_{ij,t+h}.
\end{equation}
Nhãn này là chất lượng trung bình, không phải xác suất duy trì kết nối trong
toàn bộ năm giây.

Hai lớp tổng hợp trung bình kiểu GraphSAGE~\cite{graphsage}, mỗi lớp 64 chiều
và ReLU, kết hợp thông tin nút với đặc trưng cạnh. Bộ giải mã nhận hai biểu
diễn nút và đặc trưng cạnh, qua lớp 64 chiều và sigmoid; logits được đối xứng
hóa trước sigmoid. Mô hình có 19.137 tham số. Đối chứng MLP bỏ hai nhánh truyền
thông điệp, giữ cùng đầu vào và chiều rộng lớp, có 13.505 tham số. Đây là
ablation của \emph{truyền thông điệp học được}, không phải mô hình hoàn toàn
không có thông tin topology: đặc trưng bậc nút và tỷ lệ nhận vẫn được giữ.
Số tham số hai mô hình không bằng nhau, nên chênh lệch không thể quy hoàn
toàn cho riêng cấu trúc đồ thị.
Với $\mathcal N_i$ là các hàng xóm đo được, phép cập nhật có dạng
\begin{equation}
h_i^{(\ell)}=\operatorname{ReLU}\!\left(W_s^{(\ell)}h_i^{(\ell-1)}+b_s^{(\ell)}
+W_m^{(\ell)}\left[\overline{h}_{\mathcal N_i}^{(\ell-1)}\Vert
\overline e_{\mathcal N_i}\right]+b_m^{(\ell)}\right),\quad \ell=1,2.
\end{equation}
Trung bình bằng 0 khi nút không có hàng xóm. MLP bỏ hạng chứa $W_m,b_m$;
các đặc trưng cạnh của chính cặp nút vẫn đi vào bộ giải mã. Bộ giải mã sinh
$z_{ij}$ và đầu ra là $\sigma((z_{ij}+z_{ji})/2)$, đường chéo đặt bằng 0.

Huấn luyện bằng MSE trên các cạnh không chéo, Adam với learning rate 0,001,
weight decay $10^{-5}$, batch 128, clip norm 5, tối đa 60 epoch; chọn MSE thấp
nhất trên tập chọn epoch, patience 10, seed 20261123. Hai cấu hình n5/n7 có
số mẫu cân bằng. GNN và MLP lần lượt chọn epoch 58 và 59. Suy luận triển khai
bằng NumPy trên CPU; sai số tuyệt đối lớn nhất so với PyTorch trên 40 snapshot
kiểm tra là $1{,}19\times10^{-7}$, thấp hơn ngưỡng $10^{-5}$ đã định.

\subsection{Các nhóm so sánh và phân tích}
Bốn arm gồm \code{lq-route}, \code{gnn-route}, \code{mlp-route} và
\code{olsr-etx-fast}. Ba arm đầu cùng dùng margin ETX 1,0, warmup ba bước,
dwell năm giây; chỉ thay ma trận chất lượng dùng tính chi phí tuyến. Cơ chế
phát hiện liên kết hỏng vẫn dùng tỷ lệ nhận hiện tại dưới 0,1. ETX-fast kế
thừa HELLO 0,5 s và TC 1,25 s đã chọn ở v7. Relay đứng yên trong mọi arm.

Validation gồm 80 episode: mỗi cấu hình hai khối static và tám khối gradual,
mỗi khối bốn arm. Go chỉ xét tính đầy đủ và đúng đắn kỹ thuật, không xét
dấu hoặc mức tăng PDR. Test gồm 192 episode: bốn khối static và 20 khối gradual
mỗi cấu hình. Ba bước kiểm định lần lượt là PDR của GNN trừ chất lượng hiện
tại, GNN trừ MLP, rồi GNN trừ ETX-fast; một bước chỉ được khẳng định khi cận
dưới CI 95\% dương và mọi bước trước đã đạt. Bootstrap 10.000 lần trên 40 khối
gradual ghép cặp, cân bằng n5/n7, seed 20261120. Các frame hoặc gói tin không
được coi là các mẫu độc lập.

MSE ngoại tuyến của GNN, MLP và EWMA được đo trên cùng trace test của
\code{lq-route}, tách biệt với PDR vòng kín. Thời gian tính quyết định được
collector đo quanh toàn bộ lời gọi bộ điều khiển trên server đang chạy bốn
worker; đây không phải độ trễ truyền lệnh hay thời gian tính trên UAV.
Server là máy ảo KVM với tám vCPU, báo model Intel Xeon Silver 4110
2,10 GHz. Mỗi worker giới hạn OMP, OpenBLAS và MKL về một luồng.
Chi phí báo cáo vị trí, vận tốc và RSSI chưa được mô phỏng thành lưu lượng
tranh chấp kênh, nên không áp nguyên ước lượng báo hiệu chỉ gồm tỷ lệ nhận
của v7 cho hệ thống GNN.

\section{Vòng v9: kích hoạt điều khiển vị trí bằng dự báo}\label{sec:m9}
Vòng v9 được thiết kế trước khi xem test v8, dùng nguyên trọng số GNN/MLP
đã khóa và không fit lại. Tất cả arm giữ bộ định tuyến \code{lq-route} theo
chất lượng \emph{hiện tại}; thay đổi cần kiểm tra là cơ chế kích hoạt chuyển
động. Bốn arm là \code{route-only} (đứng yên), \code{reactive-motion},
\code{gnn-motion} và \code{mlp-motion}. Cách bố trí lại relay, giới hạn tốc độ
và ngân sách giống nhau trong ba arm có chuyển động.

\subsection{Kích hoạt và luật chọn vị trí}
Chất lượng đường tốt nhất được định nghĩa bằng bài toán đường rộng nhất:
\begin{equation}
 Q_t=\max_{P:s\rightarrow d}\;\min_{(i,j)\in P}q_{ij,t}.
\end{equation}
Reactive dùng EWMA hiện tại; hai arm dự báo dùng thêm ma trận dự báo của
mình. Sau ba frame quan sát, chuyển động được kích hoạt khi $Q_t<0{,}70$
trong hai frame liên tiếp. Với GNN/MLP, điều kiện yếu có thể đến từ hiện tại
hoặc dự báo, nên vẫn giữ fallback phản ứng như reactive. Cạnh dự báo chỉ
được xét nếu tỷ lệ nhận hiện tại ít nhất 0,1.

Relay được xếp theo hình chiếu lên đường nối hai đầu mút hiện tại. Vị trí
đích là các điểm chia đều đoạn đó; không đọc vị trí đầu mút trong tương lai.
Vận tốc hướng tới vị trí đích, tối đa 8 m/s, deadband 1 m, ngân sách 400 m
cho mỗi relay. Mỗi đợt tối đa mười bước, sau đó nghỉ năm bước; dừng sớm nếu
chất lượng hiện tại và chất lượng dự báo (nếu có) đều đạt 0,80 trong ba frame.
Lệnh được kiểm tra khoảng cách nhỏ nhất liên tục trong một bước, tính cả
vận tốc đầu mút đang quan sát; nếu dưới 5 m thì không gửi chuyển động cho
các relay ở bước đó. Simulator còn thực thi ràng buộc an toàn riêng.

Để có thách thức chia cắt, hai đầu mút gradual đi xa nhau với tốc độ
\textbf{tốc độ danh định 2 m/s mỗi đầu mút} từ giây 30; các điều kiện còn lại giữ như v8.
Simulator nhân tốc độ với hệ số ngẫu nhiên theo layout seed trong [0,85; 1,15],
dùng chung cho các arm; nguồn đi theo chiều $+y$, đích theo chiều $-y$ và dừng
ở biên vùng mô phỏng. Tốc độ danh định này khác dữ liệu fit (1 m/s), vì vậy đây đồng thời là kiểm tra chuyển
miền. Chưa có dữ liệu để mặc định dự báo được hiệu chuẩn cho tốc độ mới.
Hơn nữa, dữ liệu fit chỉ có relay đứng yên, trong khi các arm v9 tạo vận tốc
relay mới. Mô hình dự báo không được huấn luyện như bộ tiên đoán phản thực
cho từng hành động; ở đây nó chỉ kích hoạt một luật bố trí đã cố định.

\subsection{Cổng khả năng kiểm tra và endpoint chính}
Validation có 80 episode như v8, dùng seed mới. Ngoài các kiểm tra kỹ thuật,
an toàn, ngân sách và relay đứng yên đúng arm, ít nhất 8/16 episode gradual
\code{route-only} phải mất kết nối ít nhất năm giây. Cổng này kiểm tra baseline
có đủ thách thức để trả lời câu hỏi; không lấy lợi ích GNN làm điều kiện go.
Nếu không đạt thì dừng test, giữ nguyên kết quả, không tăng tốc rồi chạy lại
cùng nghiên cứu.

Nếu go, test có 192 episode. Endpoint chính là số giây mất kết nối của
reactive trừ GNN; sau đó mới xét PDR GNN trừ reactive và số giây mất kết nối
của MLP trừ GNN. Quy tắc CI và thứ tự cố định giống v8. Mất kết nối được đếm
trên frame 1--240 khi nguồn và đích không nằm cùng thành phần liên thông
trong đồ thị probe ngưỡng 0,5. Chỉ số này khác số giây OLSR không có tuyến
và khác mất gói UDP. PDR tăng riêng lẻ không đủ chứng minh ngăn chia cắt.
Chia cắt ở đây là một endpoint theo ngưỡng đồ thị probe, không đồng nghĩa
mọi khả năng truyền vô tuyến giữa hai thành phần đã biến mất.
Quãng đường, số đợt chuyển động, outage dài nhất và PDR luồng nền được báo
cáo mô tả cùng kết quả chính. Quãng đường chưa được chuyển thành năng lượng.
Tất cả điều khiển vị trí mới dừng ở vận tốc trong ns-3, chưa kết nối autopilot
hoặc thử nghiệm UAV thật.
